\documentclass[12pt,a4paper]{article}

\usepackage[utf8]{inputenc}
\usepackage[T1]{fontenc}
\usepackage[margin=2.5cm]{geometry}
\usepackage{hyperref}
\usepackage{microtype}
\usepackage{parskip}
\usepackage{lmodern}

\hypersetup{colorlinks=true, urlcolor=blue!60!black}

\pagestyle{empty}

\begin{document}

\begin{flushright}
Kundai Farai Sachikonye \\
Auf der Lichtung 19, 82178 Puchheim \\
\href{https://github.com/fullscreen-triangle}{github.com/fullscreen-triangle} \\
\texttt{kundai.sachikonye@bitspark.com} \\
(+49) 1626386344 \\[6pt]
17 May 2026
\end{flushright}

\vspace{1em}

Ulrike Auer \\
EagleBurgmann Germany GmbH \& Co.\ KG \\
\texttt{Ulrike.Auer@eagleburgmann.com}

\vspace{1.5em}

\textbf{Re: Scientific Innovator --- AI-Driven Engineering (f/m/d)}

\vspace{1em}

Dear Ms Auer,

I am applying for the Scientific Innovator position. I have built original
systems at each layer the role describes --- agentic LLM workflows, surrogate
models for physics-based simulation, and a formal training principle for
domain-specific models --- and I can demonstrate all of them running live in
a browser before you read further:

\begin{itemize}
  \item Multi-model RAG pipeline with hybrid LLM + mathematical solver selection:
    \href{https://github.com/fullscreen-triangle/four-sided-triangle}{github.com/fullscreen-triangle/four-sided-triangle}
  \item Agentic OS interface (O($\log_k N$) specialist cascade routing):
    \href{https://github.com/fullscreen-triangle/zangalewa}{github.com/fullscreen-triangle/zangalewa}
  \item Domain-specific model training (Backward Training Principle):
    \href{https://github.com/fullscreen-triangle/purpose}{github.com/fullscreen-triangle/purpose}
  \item First-principles gas and fluid dynamics (Navier-Stokes from partition theory, viscosity 2.9\,\% MAE):
    \href{https://gas-giant.vercel.app/simulate}{gas-giant.vercel.app/simulate}
\end{itemize}

\textbf{Surrogate models and the surrogate training problem.}
The central challenge in building surrogate models for FEM/CFD simulation is not
the architecture --- it is the training distribution. Conventional approaches
build surrogates incrementally from less to more demanding operating conditions.
My \href{https://github.com/fullscreen-triangle/purpose}{Purpose} framework
establishes the opposite as the formally optimal strategy: the
\textbf{Backward Training Principle} proves that a model $\epsilon$-competent
on the hardest single-objective instance (the \emph{pure-intent regime}) is
$\epsilon$-competent on every less-constrained sub-regime by \emph{feasibility
inclusion}, not by generalisation. For a seal, the pure-intent regime is the
most extreme operating condition --- highest pressure, temperature, and
chemical exposure simultaneously --- from which all normal operating envelopes
are feasibility restrictions. Training the surrogate there produces a model
that works across the full design space without
re-training. A Forward Training Collapse Theorem establishes a provably
positive VC failure-gap for the conventional direction; a Sample Complexity
Theorem gives $(N+1)\times$ savings for $N$-layer constraint cascades.
I have already applied this pattern in practice:
\href{https://github.com/fullscreen-triangle/homo-veloce}{homo-veloce}
is a surrogate model system that distils physics-based biomechanics solvers
(Gaussian processes and graph neural networks over physical simulation outputs)
into lightweight neural networks for real-time inference --- the same pattern
FEM/CFD surrogate development requires.

\textbf{CFD and fluid simulation from first principles.}
The position lists FEM/CFD simulation as one of the skill areas and surrogate
modelling as the target application. The
\href{https://gas-giant.vercel.app/simulate}{Gas Giant} platform I have built
takes a different approach: rather than training a surrogate on CFD data, it
derives the governing equations from the same bounded phase space axiom.
Kirchhoff circuit laws map analytically to the Navier-Stokes equations;
viscosity emerges as $\mu = \tau_c \times g$ (partition lag times coupling
strength) rather than as a phenomenological fit. The gas-to-liquid transition
occurs at network density thresholds $\rho_C \approx 0.3$--$0.7$, with no
fitted transition parameter.
Validation: 2.9\,\% mean absolute error across 12 pure liquids for viscosity;
0.9\,\% mean error for Poiseuille flow parabolic velocity profiles.
Temperature emerges as $T = (\hbar/k_B) \cdot dM/dt$ --- a processing
rate identity --- from which the ideal gas law $PV = Nk_BT$ and the
adiabatic invariant $PV^\gamma = \mathrm{const}$ follow as corollaries,
validated to $U/({\tfrac{3}{2}}Nk_BT) = 1.000 \pm 0.2\%$.
For sealing technology specifically, the fluid dynamics around the seal face ---
viscosity, diffusion, pressure gradients, flow regime transitions --- are
exactly what the /fluids tool computes, from partition geometry rather than
from empirical correlations.

\textbf{Digital twin and condition monitoring.}
The \href{https://github.com/fullscreen-triangle/syndrome}{Syndrome} framework
encodes system state as a vector $\mathbf{D} \in [0,1]^8$ across
eight monitoring dimensions, each with a coherence index $\eta_i \in [0,1]$
computed from sensor time-series regularity. The framework computes
$\mathcal{O}(N \log N)$ state trajectories and selects the
minimum-intervention maintenance action via sparse $\ell_1$ linear programming.
This architecture was developed for biological systems but the formalism is
domain-agnostic: the eight oscillator classes map directly onto the relevant
condition monitoring dimensions for rotating or dynamic sealing equipment
(leakage rate, vibration spectrum, temperature distribution, pressure differential,
wear pattern, corrosion state, fatigue index, and interfacial geometry).
The sparse LP answers the predictive maintenance question directly: which
component to address, when, and in what order, to minimise downtime and
avoid cascading failure.
The \href{https://github.com/fullscreen-triangle/verum}{Verum} framework
adds a complementary state-reconstruction layer: a vehicle modelled as a
20-node oscillatory circuit graph, where each subsystem's state is its
categorical depth $V_i = -\log_2 P(\sigma_i)$ and edge conductances derive
from a universal transport formula unifying electrical, thermal, viscous, and
diffusive coupling. A Banach contraction mapping (Lipschitz constant
$\lambda \in [0.3, 0.7]$) guarantees reconstruction of the complete state from
33\,\% surface observations alone --- validated on 2023 Formula One telemetry
(Bahrain GP, 13,794 samples, convergence residual $7.2 \times 10^{-9}$,
ICE-Turbo fault detected 13~laps before failure event). The same architecture
maps directly onto sealing equipment: the circuit graph nodes are pressure
differential, temperature, vibration, wear state, and interfacial geometry;
the surface sensors are the observed boundary conditions; the complete
condition monitoring vector is the hidden state recovered by the contraction.

\textbf{Agentic workflows and LLM engineering assistants.}
\href{https://github.com/fullscreen-triangle/four-sided-triangle}{Four-Sided-Triangle}
is an 8-stage metacognitive RAG pipeline with adaptive hybrid selection between
LLM reasoning and mathematical solvers, backed by a Rust engine providing
10--50$\times$ speedup over Python for compute-intensive stages.
\href{https://github.com/fullscreen-triangle/zangalewa}{Zangalewa} is the
Minimum Sufficient Interceptor sitting at the boundary of a blank-screen OS:
a user types a natural language request; the system extracts semantic coordinates
via a hierarchical $k$-ary specialist cascade requiring $O(\log_k N)$
small-model invocations, and routes to the correct tool or expert sub-model.
A Sufficiency Theorem proves the parameter budget is $O(D \cdot b \cdot \log|\Sigma|)$,
independent of the downstream tool space size. For an LLM engineering
assistant that must route between FEM queries, material databases, maintenance
logs, and design automation tools, this is the architecture that keeps the
routing layer small and the specialist tools deep.

\textbf{Distributed sensor coordination.}
A digital twin for sealing technology aggregates sensor streams from multiple
physical sites under real-time timing and security constraints. The
\href{https://github.com/fullscreen-triangle/pylon}{Pylon} framework I have
developed addresses this directly: it proves a structural identity between
distributed networks and ideal gases ($PV = Nk_BT$, validated to 0.6\%
on a 128-node, $10^6$-packet testbed) and derives coordination protocols
from that identity, achieving $33\times$ throughput improvement, $20\times$
jitter reduction, and $1000\times$ faster recovery versus conventional
approaches. Security derives from the Second Law of Thermodynamics rather
than cryptographic hardness: any coordinated sensor spoofing attack requires
a thermodynamically forbidden entropy reduction, giving 100\% detection with
0\% false positives. The precision-by-difference synchronisation layer
($\Delta P(v_i, k) = T_\mathrm{ref}(k) - t_i(k)$ against an atomic clock
reference) provides 87\,ns phase coherence across nodes --- the timing
substrate that makes the condition-monitoring state vector $\mathbf{D}$
in the Syndrome framework trustworthy at the sensor level.

\textbf{Background.}
I hold an MSc in Computational Biology (Universität Tübingen), an MSc in
Pharmaceutical Biotechnology (HAW Hamburg), and completed doctoral research
in computational science at TUM. My primary languages are Python (PyTorch,
scikit-learn, Ray), Rust for performance-critical components, and TypeScript
for deployment. I am legally resident in Germany with full work authorisation,
available immediately, and the Wolfratshausen location is a short commute
from where I currently live. English is native; German is conversational.

\vspace{1.5em}
Yours sincerely,

\vspace{2.5em}
\textbf{Kundai Farai Sachikonye}

\end{document}
